% Nội dung phần 5: Hạn chế và hướng phát triển.

\section{Hạn chế và hướng phát triển}

\subsection{Hạn chế của phân tích}

\subsubsection{Hạn chế về thiết kế nghiên cứu}

\begin{itemize}
  \item \textbf{Không có treatment/control ngẫu nhiên.} So sánh hai nhóm giá là
  phân tích quan sát. Balance diagnostics (Hình~\ref{fig:ab-balance}) cho thấy
  nhiều covariate có $|\text{SMD}| > 0{,}1$, tức hai nhóm khác nhau một cách hệ
  thống ngay từ đầu. Việc điều chỉnh bằng mô hình chỉ xử lý được các yếu tố đã
  quan sát; residual confounding bởi đặc điểm sản phẩm, danh mục và thói quen
  duyệt vẫn có thể tồn tại.

  \item \textbf{Giả thuyết được định hướng từ EDA.} Nhóm giá trở thành đối tượng
  kiểm định sau khi quan sát dữ liệu, không phải trước đó. Vì vậy p-value trong
  báo cáo cần được đọc cùng effect size, balance diagnostics, kiểm tra ổn định
  theo thời gian và hiệu chỉnh Holm, chứ không đứng một mình.

  \item \textbf{Quan hệ theo \texttt{Order} chịu cấu trúc chọn lọc.} Tại mỗi vị
  trí $t$, chỉ những session còn tồn tại đến $t$ được quan sát. Xu hướng giảm của
  xác suất kết thúc theo tiến trình vì vậy trộn lẫn hai cơ chế: thay đổi hành vi
  thực sự và việc các session có xu hướng dừng sớm đã bị loại khỏi mẫu.
\end{itemize}

\subsubsection{Hạn chế về dữ liệu}

\begin{itemize}
  \item \textbf{Không có thông tin chuyển đổi.} Dữ liệu không ghi giao dịch, giỏ
  hàng hay thanh toán. Vì vậy \texttt{Exit} chỉ có nghĩa ``click cuối của
  session'', hoàn toàn không phân biệt được một phiên kết thúc vì mua hàng thành
  công với một phiên kết thúc vì thất vọng bỏ đi. Đây là hạn chế nghiêm trọng
  nhất đối với giá trị kinh doanh của mô hình.

  \item \textbf{Không có chiều thời gian trong phiên.} Không có giờ click và
  không có thời lượng giữa hai click liên tiếp -- những tín hiệu thường mạnh nhất
  trong bài toán dự báo rời bỏ. Mô hình vì vậy buộc phải thay thế bằng các đại
  lượng đếm (số sản phẩm, số danh mục đã xem).

  \item \textbf{Không có định danh người dùng xuyên phiên.} Không thể phân biệt
  khách quay lại với khách mới, cũng không thể xây dựng đặc trưng lịch sử dài hạn.

  \item \textbf{Phụ thuộc trong cụm.} Các click trong cùng một session không độc
  lập. Nhóm đã dùng cluster-robust covariance và cluster bootstrap theo session,
  nhưng đây chỉ là hiệu chỉnh cho suy luận, không mô hình hóa tường minh cấu trúc
  phụ thuộc; một mô hình hiệu ứng hỗn hợp hoặc survival có thể phù hợp hơn.

  \item \textbf{Tháng 8 không đầy đủ.} Dữ liệu dừng ở ngày 13/08 và riêng ngày
  này chỉ có 200 click nên đã bị loại. Temporal test vì vậy chỉ gồm 12 ngày với
  13.948 click, khiến khoảng tin cậy bootstrap khá rộng
  (Bảng~\ref{tab:final-metrics}).

  \item \textbf{Tính đại diện.} Dữ liệu từ năm 2008 và của một cửa hàng cụ thể ở
  Ba Lan. Giao diện thương mại điện tử, thiết bị truy cập và hành vi người dùng
  đã thay đổi rất nhiều; cần đánh giá lại toàn bộ trước khi áp dụng cho môi
  trường hiện tại.
\end{itemize}

\subsubsection{Hạn chế về mô hình}

\begin{itemize}
  \item \textbf{Concurvity cao trong GAM.} Mức concurvity lớn nhất giữa các term
  đạt 0,994 (Bảng~\ref{tab:collinearity}). Đây là lý do các partial-effect plot ở
  Hình~\ref{fig:gam-smooths} chỉ được dùng để mô tả dạng hàm, không được diễn
  giải như hiệu ứng riêng của từng biến.

  \item \textbf{Biến \texttt{product\_model} bị bỏ qua.} Với 217 mức, biến này
  không được đưa vào mô hình chính vì chưa có quy tắc gộp mức hiếm được kiểm
  chứng. Có thể vì vậy mà một phần thông tin về sản phẩm bị bỏ sót.

  \item \textbf{Khả năng phân biệt còn khiêm tốn.} ROC--AUC 0,666 và Brier Skill
  Score 0,039 chưa đủ để coi mô hình là một hệ thống sẵn sàng triển khai.
\end{itemize}

\subsection{Đề xuất một thí nghiệm A/B ngẫu nhiên thực sự}
\label{subsec:future-ab}

Để chuyển từ kết luận liên hệ sang kết luận nhân quả, cần một thí nghiệm được
thiết kế đúng. Nhóm đề xuất thiết kế sau.

\textbf{Đơn vị ngẫu nhiên hóa.} Randomize ở \emph{cấp session}, lưu
\texttt{experiment\_id} cho từng phiên và đảm bảo mỗi session chỉ thuộc đúng một
nhánh trong suốt vòng đời của nó. Ngẫu nhiên hóa ở cấp click sẽ vi phạm giả định
độc lập vì các click trong cùng phiên phụ thuộc nhau.

\textbf{Treatment.} Thay vì tùy tiện thay đổi giá sản phẩm -- điều vừa rủi ro về
doanh thu vừa khó hoàn tác -- treatment khả thi hơn là cách \emph{trình bày}
thông tin giá:

\begin{itemize}
  \item \textbf{Control:} giao diện hiển thị giá hiện tại.
  \item \textbf{Treatment:} bổ sung thông điệp giá trị hoặc ưu đãi vận chuyển,
  giữ nguyên giá sản phẩm.
\end{itemize}

\textbf{Outcome.} Primary outcome là tỷ lệ session kết thúc trong một cửa sổ
click được xác định trước (ví dụ trong vòng 3 click kể từ khi nhìn thấy
treatment). Secondary outcomes gồm độ dài session và số sản phẩm duy nhất đã xem.
Cửa sổ outcome phải được khóa \emph{trước} khi thu thập dữ liệu.

\textbf{Phân tích.} Intention-to-treat, kiểm định hai phía tại
$\alpha = 0{,}05$, không dừng sớm dựa trên p-value chưa điều chỉnh. Nếu cần theo
dõi giữa chừng, phải dùng thiết kế sequential có kiểm soát sai lầm loại I.

\textbf{Về cỡ mẫu.} Nhóm cố ý \emph{không} báo cáo một con số cỡ mẫu tính từ hai
tỷ lệ click quan sát ở Bảng~\ref{tab:ab-rate}. Lý do là thí nghiệm tương lai
randomize theo session và có estimand hoàn toàn khác (tỷ lệ ở cấp session, không
phải cấp click), nên đưa ra con số như vậy sẽ gây hiểu nhầm. Trước khi tính
power cần khóa bốn tham số: cửa sổ outcome, tỷ lệ nền ở cấp session, minimum
detectable effect có ý nghĩa kinh doanh, và mức hao hụt dự kiến. Cỡ mẫu sau đó
được tính trực tiếp theo số session mỗi nhánh.

\subsection{Hướng phát triển khác}

\begin{itemize}
  \item \textbf{Mô hình sinh tồn thời gian rời rạc.} Bài toán hiện tại về bản
  chất là ước lượng hazard rời rạc. Một mô hình survival có thể xử lý tường minh
  cấu trúc chọn lọc theo \texttt{Order} thay vì chỉ điều chỉnh gián tiếp.

  \item \textbf{Mô hình hiệu ứng hỗn hợp.} Thêm random intercept theo session sẽ
  mô hình hóa trực tiếp phần phụ thuộc trong cụm thay vì chỉ hiệu chỉnh sai số
  chuẩn.

  \item \textbf{Đối chứng bằng mô hình phi tham số.} Gradient boosting hoặc
  random forest có thể cho biết trần hiệu năng thực sự của bộ đặc trưng hiện tại;
  nếu chúng cũng chỉ đạt ROC--AUC quanh 0,67 thì giới hạn nằm ở dữ liệu chứ không
  ở dạng mô hình.

  \item \textbf{Gộp mức của \texttt{product\_model}.} Xây dựng quy tắc gộp 217 mã
  sản phẩm dựa trên tần suất trong tập train, rồi đánh giá phần thông tin bổ sung
  bằng đúng pipeline hiện tại.
\end{itemize}
